\section{Comparative analysis of members and casual riders}
\label{sec:analysis}

Unless stated otherwise, all percentages in this section are computed on each group's rides (3,897,958 member rides, 2,114,500 casual rides). In ``A vs B'' comparisons, A refers to casual riders and B to members.

\subsection{Volume and duration}

\begin{table}[!htb]
\caption{Summary by rider type}
\label{tab:riders}
\small
% Generated by scripts/generate_latex_tables.py - do not edit by hand
\begin{tabularx}{\linewidth}{Y r r r}
\toprule
\thead{Indicator} & \thead{\textcolor{member}{Members}} & \thead{\textcolor{casual}{Casual riders}} & \thead{All} \\
\midrule
Rides & 3,897,958 & 2,114,500 & 6,012,458 \\
Share of rides & 64.8\% & 35.2\% & 100\% \\
Median duration (min) & 8.69 & 11.17 & 9.46 \\
Mean duration (min) & 12.60 & 21.11 & 15.59 \\
1st quartile / 3rd quartile (min) & 5.19 / 14.63 & 6.49 / 20.04 & 5.58 / 16.33 \\
90th percentile (min) & 23.42 & 36.05 & 27.29 \\
Rides of at least 20 min & 14.1\% & 25.1\% & 18.0\% \\
Rides of at least 30 min & 5.4\% & 13.5\% & 8.3\% \\
Rides longer than 24 hours & 1,108 & 4,196 & 5,304 \\
\bottomrule
\end{tabularx}

\par\smallskip{\footnotesize\color{muted}Source: \file{results/analysis/a01\_usagers.csv}, \file{a02\_durees\_tranches.csv}. Durations in minutes, computed from the raw timestamps.}
\end{table}

\Mes{} Nearly two rides out of three are taken by members. Casual rides are longer across the whole distribution: a median 29\% higher, a 90th percentile 54\% higher, and 13.5\% of rides lasting at least 30~min vs 5.4\%. The gap depends strongly on the season: in January, the medians are almost identical (7.13 vs 7.17~min); in June, they are 12.35 vs 9.47~min.

\subsection{Day of the week}

\begin{table}[!htb]
\caption{Distribution of rides and median duration by day of departure}
\label{tab:days}
\small
% Generated by scripts/generate_latex_tables.py - do not edit by hand
\begin{tabularx}{\linewidth}{Y r r r r}
\toprule
 & \multicolumn{2}{c}{\thead{Share of the group's rides}} & \multicolumn{2}{c}{\thead{Median duration (min)}} \\
\cmidrule(lr){2-3}\cmidrule(lr){4-5}
\thead{Day of departure} & \thead{Members} & \thead{Casual} & \thead{Members} & \thead{Casual} \\
\midrule
Monday & 14.0\% & 11.5\% & 8.41 & 10.83 \\
Tuesday & 16.1\% & 11.4\% & 8.50 & 9.89 \\
Wednesday & 16.2\% & 11.7\% & 8.49 & 9.68 \\
Thursday & 15.8\% & 12.5\% & 8.51 & 9.93 \\
Friday & 14.7\% & 15.7\% & 8.57 & 11.00 \\
Saturday & 12.8\% & 21.2\% & 9.55 & 13.06 \\
Sunday & 10.5\% & 16.0\% & 9.23 & 12.82 \\
\midrule
\textbf{Weekend (Sat.\,+\,Sun.)} & \textbf{23.3\%} & \textbf{37.2\%} & & \\
\bottomrule
\end{tabularx}

\par\smallskip{\footnotesize\color{muted}Source: \file{a03\_jour.csv}. A uniform distribution would put 28.6\% of rides at the weekend.}
\end{table}

\Mes{} The weekend accounts for 37.2\% of casual rides vs 23.3\% of member rides; Saturday is the busiest day for casual riders (21.2\%). Members ride mostly from Tuesday to Thursday. The gap in median duration is largest on Saturday (13.06 vs 9.55~min).

\subsection{Month and seasonality}

\begin{table}[!htb]
\caption{Rides and median duration by month of departure}
\label{tab:months}
\small
% Generated by scripts/generate_latex_tables.py - do not edit by hand
\begin{tabularx}{\linewidth}{L{1.9cm} R R R R R}
\toprule
\thead{Month of departure} & \thead{Member rides} & \thead{Casual rides} & \thead{Casual share} & \thead{Median, members (min)} & \thead{Median, casual (min)} \\
\midrule
Oct. 2025 & 414,231 & 214,849 & 34.2\% & 8.57 & 10.74 \\
Nov. 2025 & 252,033 & 94,930 & 27.4\% & 7.92 & 8.97 \\
Dec. 2025 & 109,411 & 27,173 & 19.9\% & 7.49 & 7.69 \\
Jan. 2026 & 109,974 & 23,962 & 17.9\% & 7.17 & 7.13 \\
Feb. 2026 & 157,270 & 40,162 & 20.3\% & 7.43 & 8.21 \\
Mar. 2026 & 223,698 & 85,031 & 27.5\% & 7.87 & 9.88 \\
Apr. 2026 & 308,523 & 127,309 & 29.2\% & 8.20 & 10.05 \\
May 2026 & 399,616 & 236,362 & 37.2\% & 8.87 & 11.83 \\
June 2026 & 444,235 & 297,844 & 40.1\% & 9.47 & 12.35 \\
July 2026 & 500,577 & 346,209 & 40.9\% & 9.36 & 11.98 \\
Aug. 2026 & 511,153 & 336,005 & 39.7\% & 9.33 & 11.76 \\
Sept. 2026 & 467,237 & 284,664 & 37.9\% & 8.90 & 11.08 \\
\midrule
\textbf{Total} & \textbf{3,897,958} & \textbf{2,114,500} & \textbf{35.2\%} & & \\
\bottomrule
\end{tabularx}

\par\smallskip{\footnotesize\color{muted}Source: \file{a04\_mois.csv}. March 2026: possible gap for rides crossing midnight on the 31st (Section~\ref{sec:data}).}
\end{table}

\Mes{} Both groups follow the same calendar, with very different amplitudes: ×14.4 between January and July for casual riders, ×4.7 between December and August for members. From May to September, casual riders take 71.0\% of their annual rides (59.6\% for members); their share of all rides ranges from 17.9\% in January to 40.9\% in July.

\subsection{Hour of departure}

\begin{table}[!htb]
\caption{Hourly indicators}
\label{tab:hours}
\small
% Generated by scripts/generate_latex_tables.py - do not edit by hand
\begin{tabularx}{\linewidth}{Y r r}
\toprule
\thead{Hourly indicator (local time of departure)} & \thead{Members} & \thead{Casual} \\
\midrule
Departures 7--9 a.m., \% of weekday rides & 15.1\% & 7.6\% \\
Departures 5--6 p.m., \% of weekday rides & 11.9\% & 10.8\% \\
Departures 10 a.m.--4 p.m., \% of weekday rides & 26.3\% & 30.7\% \\
Departures 10 a.m.--4 p.m., \% of weekend rides & 44.6\% & 45.0\% \\
Weekday peak hours (7--9 a.m. and 4--7 p.m.), \% of all the group's rides & 35.3\% & 23.1\% \\
\bottomrule
\end{tabularx}

\par\smallskip{\footnotesize\color{muted}Source: \file{a05\_heure.csv}. Local time of departure; ``7--9 a.m.'' covers departures from 7:00 to 8:59 a.m.}
\end{table}

\Mes{} On weekdays, the morning peak is twice as pronounced among members (15.1\% vs 7.6\%), while both groups peak at 5 p.m. At weekends, their hourly profiles are almost identical. \Int{} The members' double peak is consistent with commuting; part of the casual weekday rides follows the same rhythm.

\subsection{Bike type and interactions}

\begin{table}[!htb]
\caption{Bike type and type of day}
\label{tab:bikes}
\small
% Generated by scripts/generate_latex_tables.py - do not edit by hand
\begin{tabularx}{\linewidth}{Y r r r r}
\toprule
 & \multicolumn{2}{c}{\thead{Share of the group's rides}} & \multicolumn{2}{c}{\thead{Median duration (min)}} \\
\cmidrule(lr){2-3}\cmidrule(lr){4-5}
\thead{Bike, type of day} & \thead{Members} & \thead{Casual} & \thead{Members} & \thead{Casual} \\
\midrule
Classic, weekday & 24.2\% & 15.3\% & 8.92 & 14.32 \\
Classic, weekend & 7.5\% & 10.7\% & 10.09 & 17.79 \\
Electric, weekday & 52.5\% & 47.5\% & 8.33 & 9.39 \\
Electric, weekend & 15.8\% & 26.5\% & 9.13 & 11.55 \\
\midrule
\textbf{E-bike (total)} & \textbf{68.3\%} & \textbf{74.1\%} & & \\
\bottomrule
\end{tabularx}

\par\smallskip{\footnotesize\color{muted}Source: \file{a06\_velo.csv}, \file{a07\_interactions.csv}. The excluded rides shorter than 60~s were all e-bike rides.}
\end{table}

\Mes{} E-bikes account for the majority of rides in both groups. The duration gap is largest for classic bikes at weekends (17.79 vs 10.09~min) and smallest for e-bikes on weekdays (9.39 vs 8.33~min), the most frequent combination in both groups.

\subsection{Round trips}

\begin{table}[!htb]
\caption{Rides ending at their start station}
\label{tab:roundtrips}
\small
% Generated by scripts/generate_latex_tables.py - do not edit by hand
\begin{tabularx}{\linewidth}{Y R R R R}
\toprule
\thead{Group} & \thead{Rides with both stations} & \thead{Round trips} & \thead{Share of round trips} & \thead{Median round trip (min)} \\
\midrule
Members & 2,652,984 (68\%) & 58,539 & 2.2\% & 14.38 \\
Casual riders & 1,396,510 (66\%) & 103,487 & 7.4\% & 16.70 \\
\bottomrule
\end{tabularx}

\par\smallskip{\footnotesize\color{muted}Source: \file{a08\_boucles.csv}. Computed only for rides with both stations; the share of rides covered is shown in parentheses.}
\end{table}

\Mes{} Round trips are 3.4 times more frequent among casual riders. \Int{} This points to leisure use, to be confirmed by a station-level analysis.
